\section{Introduction}\label{sec:introduction}

The Virasoro conjecture organizes the descendant Gromov--Witten invariants of a smooth projective variety into a system of differential equations. Its simplest instance is the intersection theory of the moduli spaces of curves: Witten's conjecture and Kontsevich's theorem identify the corresponding generating function with the distinguished solution of the KdV hierarchy \cite{Witten1991,Kontsevich1992}. Eguchi, Hori, and Xiong proposed a target-dependent Virasoro system for Fano varieties \cite{EHX1997}. Eguchi, Jinzenji, and Xiong extended the free-field formulation to odd and off-diagonal Hodge classes, crediting Katz's proposal to use a Hodge index in the grading \cite{EJX1998}. We use the first-Hodge-degree superspace operators and central normalization of Getzler \cite{Getzler1999}.

For a smooth projective variety $X$, let $F_g^X(t)$ be the generating series of its ordinary, unreduced genus-$g$ descendant invariants. The variables $t_m^a$ record insertions $\tau_m(\phi_a)$ for a homogeneous cohomology basis $\{\phi_a\}$, and Novikov variables record curve classes. The total descendant potential is
\[
 Z_X(t)=\exp\!\left(\sum_{g\geq0}\h^{g-1}F_g^X(t)\right).
\]
The conjecture asserts
\[
 L_k^X Z_X=0\qquad(k\geq-1),
\]
where the target-dependent differential operators $L_k^X$ are formed from the Poincar\'e pairing, the first Hodge index $p$ on $H^{p,q}(X)$, and cup product with $c_1(TX)$. Their precise quantization, super signs, and scalar normalization are fixed in \cref{sec:conventions}. All equations are formal coefficientwise identities.

\begin{theorem}\label{thm:main}
Let \(X\subset\PP^N_\C\) be a smooth complete intersection. Its ordinary, unreduced total descendant potential satisfies the Virasoro constraints in every genus and curve class, with arbitrary insertions from \(H^*(X;\C)\).
\end{theorem}

Thus \cref{thm:main} resolves the Virasoro conjecture positively for projective-space complete intersections. Primitive and odd insertions are included, with no semisimplicity hypothesis.

The genus-zero constraints are known in general. Liu and Tian proved the decisive $L_1$ and $L_2$ identities for even-class inputs without a Fano hypothesis \cite[Theorem~0.2]{LiuTian1998}. Dubrovin and Zhang proved the genus-zero Frobenius-manifold formulation and the semisimple genus-one case \cite[Main Theorem]{DubrovinZhang1999}; Getzler treated the genus-zero superspace formulation used here \cite[Section~4]{Getzler1999}.

Several all-genus cases are also established. Getzler proved the ordinary constraints for Calabi--Yau targets with $c_1(X)=0$ and $H^1(X)=0$ \cite[Theorem~7.1]{Getzler1999}. Okounkov and Pandharipande proved the fixed-degree relative constraints for target curves, including odd descendants \cite[Theorem~3]{OP2003}. Givental's quantized semisimple construction, together with Teleman's classification, covers the corresponding homogeneous semisimple geometric theories \cite{Givental2001,Teleman2012}. More recently, Guo, Zhang, and Zhou proved genus-one constraints on the ambient state space of smooth Fano complete intersections and announced further full-cohomology genus-one extensions for several families, with details deferred to a forthcoming paper \cite[Theorem~1 and Remark~3.10]{GuoZhangZhou2026}. The present argument addresses all genera and all insertions simultaneously through degeneration.

\subsection{Degeneration and primitive cohomology}

Degeneration expresses Gromov--Witten invariants through relative invariants of simpler pieces \cite{Li2001,Li2002,AF2016}. Relative reconstruction and equation splitting were developed by Maulik and Pandharipande \cite{MaulikPandharipande2006}. Arg\"uz, Bousseau, Pandharipande, and Zvonkine used nodal invariants to reconstruct the full Gromov--Witten theory of complete intersections by induction on dimension and defining degrees \cite[Theorem~5.3 and Section~5.3]{ABPZ}. We use the same geometric reduction; the additional objective is to transport the Virasoro equations, including their first-Hodge-weighted contractions.

The reduction has two stages. Factoring one defining equation and resolving the resulting family gives a degeneration
\[
 X\rightsquigarrow U\cup_D\widetilde M,
 \qquad \widetilde M=\Bl_S M.
\]
Here $U$ and $M$ are complete intersections of smaller defining degree, while the seam $D$ and blowup center $S$ have smaller dimension. We first prove the constraints for $\widetilde M$ from those for $M$ and projective-bundle towers over $S$. We then prove them for $X$ from those for $U$, $\widetilde M$, and a ruled bundle over $D$. The needed bundles are towers of projectivizations of sums of line bundles, so toric-bundle transfer supplies their constraints from those of their bases \cite{CGT}. \Cref{ind:section} carries out this induction and verifies the first-Hodge convention in the bundle transfer.

Two difficulties prevent an immediate deduction from the numerical degeneration formula. Primitive insertions need not extend individually through a degeneration. Moreover, the Virasoro operators weight contracted indices by their first Hodge degrees, so their contractions cannot be treated as ungraded diagonal insertions. For fixed genus, degree, and descendant orders, we regard the coefficient of $L_k^X Z_X$ as a multilinear error tensor in its cohomology inputs. We transport scalar contractions of this tensor with coherently extending classes of fixed Hodge type, graded inverse pairings, and ordinary Gromov--Witten tensors from independent copies of the theory. These are the tests that will detect the error without extending each primitive input.

\subsection{The transport argument}

The transport needed in both stages rests on three constructions.

First, contractions are performed on the first page of the weight spectral sequence. Its pairing is perfect at the level of complexes and has an inverse tensor supported at vertices and edges of the degeneration graph. The first Hodge degree, including Tate twists, commutes with its differential. Multiplicative comparison and trace then evaluate the contracted tensor without choosing representatives in the surviving cohomology. We use Steenbrink's limiting mixed Hodge structures and Fujisawa's ordered model, multiplicative comparison, and trace \cite{Steenbrink1976,Steenbrink1977,Fujisawa2008,Fujisawa2014}.

For this calculation, the virtual class pushed forward by evaluation at the recorded markings must extend to a resolved product of the family. Configurations of points on common expansions provide that product \cite{AFConfigurations}. We compute its component restrictions by virtual splitting before integration, keeping the recorded markings' component assignments. When disconnected groups of maps are separated, we retain the joint evaluation of the group carrying those markings. These correspondence statements are stronger than a numerical degeneration identity and are proved in \cref{ev:section}.

Second, relative caps reproduce the contact pairing at a seam. Ordinary descendants on one cap realize functionals on a bounded incoming contact space; descendants on the other prepare the corresponding states with all unwanted outgoing profiles canceled. The nonzero projective-space fiber invariants of Hu, Li, and Ruan \cite{HLR2008} give the invertible linearization from which the full disconnected cap construction begins. In the blowup configuration, corrections of counting degree zero strictly decrease contact. Together with degree bounds, this makes every required inversion coefficientwise finite.

Insert many widely separated switches into a degeneration with a long chain of ruled components. At a switch, either keep the original joining or separate the two sides and add the prepared caps. The cap tests may use many auxiliary markings, but their ordinary insertions are integrated into the evaluation correspondence. A quadratic Virasoro operator modifies at most two existing labels. Thus the exposed cohomological operations occupy a bounded number of positions, independently of the auxiliary markings. The alternating sum over switch choices cancels by toggling the first switch away from these positions; see \cref{fig:sewing}. Every nonempty surgery subset smooths to a disjoint union of the shorter targets described above, whose absolute constraints are known. The cancellation therefore proves every such scalar test of the error on the original target.

Third, these tests detect all primitive tensors. Hodge--Riemann gives a positive Hermitian form on primitive middle cohomology. Pairing a Virasoro error tensor with its conjugate in a second replica produces its squared norm. The contraction is among the allowed tests, and positivity forces the tensor itself to vanish. This argument has no bound on the number of primitive insertions.

Transport initially gives coefficients selected by the divisor degrees that extend coherently through the degenerations. To obtain the asserted result in each curve class, we verify that these degrees separate all relevant classes, retaining additional divisor degrees for exceptional surfaces and blowups. Deformation invariance and the same scalar-test argument then give the constraints on every smooth member. These are part of the final induction, not extra hypotheses on \cref{thm:main}.

\begin{figure}[ht]
\centering
\begin{tikzpicture}[x=1cm,y=1cm,>=Latex,
  piece/.style={draw,rounded corners=2pt,minimum width=.68cm,minimum height=.55cm,fill=blue!4},
  cappiece/.style={draw,rounded corners=2pt,minimum width=.68cm,minimum height=.55cm,fill=green!7},
  busy/.style={draw=orange!75!black,fill=orange!9,rounded corners=3pt}]
\filldraw[busy] (1.35,.35) rectangle (3.15,1.25);
\filldraw[busy] (8.2,.35) rectangle (10,1.25);
\node[font=\scriptsize,anchor=south] at (2.25,1.30) {operation positions};
\node[font=\scriptsize,anchor=south] at (9.1,1.30) {operation positions};
\foreach \i/\x/\s in {0/0/U,1/1.8/C,2/3.55/C,3/6.85/C,4/8.6/C,5/10.4/V}
  {\node[piece] (t\i) at (\x,.65) {\(\s\)};}
\draw (t0)--(t1);
\draw (t1)--(t2);
\draw (t2)--(t3);
\draw (t3)--(t4);
\draw (t4)--(t5);
\draw[densely dashed,thick] (5.2,.25)--(5.2,1.07);
\node[font=\scriptsize,above] at (5.2,1.07) {idle switch};
\draw[->,thick] (5.2,-.05)--(5.2,-.65);
\node[font=\small,right] at (5.4,-.37) {toggle};
\foreach \i/\x/\s in {0/0/U,1/1.8/C,2/3.55/C,3/6.85/C,4/8.6/C,5/10.4/V}
  {\node[piece] (b\i) at (\x,-1.25) {\(\s\)};}
\node[cappiece] (Q) at (4.6,-1.25) {\(Q\)};
\node[cappiece] (P) at (5.8,-1.25) {\(P\)};
\draw (b0)--(b1);
\draw (b1)--(b2);
\draw (b2)--(Q);
\draw (P)--(b3);
\draw (b3)--(b4);
\draw (b4)--(b5);
\draw[decorate,decoration={brace,mirror,amplitude=4pt}] (4.2,-1.65)--(6.2,-1.65);
\node[font=\scriptsize,below] at (5.2,-1.82) {identity cap tests};
\end{tikzpicture}
\par

\caption{One idle switch in the transport argument. Here \(U,V\) are the ends, \(C\) is a ruled neck, and \(P,Q\) are replacement caps; see \cref{sew:configurations}. The marked neighborhoods carry the fixed cohomological operations. At a switch outside those neighborhoods, the prepared combination of cap tests reproduces the original contact identity. Toggling that switch changes the inclusion--exclusion sign and preserves the retained tensor data.}
\label{fig:sewing}
\end{figure}

\subsection{Organization and scope of the inputs}

\Cref{sec:conventions} specifies the operators, tensor tests, and completions. \Cref{fp:section,ev:section} establish the local Hodge calculation and its evaluation correspondence. \Cref{cap:section} constructs the cap tests, and \cref{sew:section} proves their transport theorem. \Cref{sec:detection} gives the positivity detector. \Cref{ind:section} assembles these results in the complete-intersection induction, including the surface and curve-class arguments.

Established degeneration, limiting-Hodge, toric-bundle, and intersection-theoretic results are used with their stated hypotheses. The local correspondence calculation, bounded cap preparation, grouped switch cancellation, and primitive tensor detector are proved here. In particular, the sewing theorem is not attributed to the numerical degeneration formula. The toric-bundle step explicitly passes from a half-total-degree convention to first Hodge degree at Hodge-neutral auxiliary parameters, with the corresponding central normalization.
